% !TEX program = LuaLaTeX+se
% !TEX encoding = UTF-8 Unicode

\documentclass[oneside, 12pt, a4paper]{book}
%\usepackage[letterpaper,showframe]{geometry}
\usepackage{fullpage} % redukcja marginesów
\usepackage{geometry}
\geometry{
	left=12mm,
	right=12mm,
	top=12mm,
	bottom=12mm,
	marginparwidth=0mm,
	marginparsep=0mm
	}
\usepackage{gregoriotex}
\gresetlinecolor{gregoriocolor}
\gresetheadercapture{annotation}{greannotation}{string}
% Format annotation above initial
\grechangestyle{annotation}{\small\bfseries}%
\grechangedim{intersyllablespacenotes}{0.18cm}{scalable}%

\usepackage{fullpage}
\usepackage{multicol}
\usepackage{setspace} % odstępy między liniami

%  Konfiguracja kolumn
\usepackage{paracol}
\setcolumnwidth{14cm,4cm}
\setlength{\columnsep}{0.4cm}
% \setlength{\columnseprule}{0.4pt}
% \colseprulecolor{red}

\usepackage{tcolorbox}
\newtcolorbox{nutybox}{size=minimal, colback=white, colframe=red, arc=0mm, width=\linewidth, halign=left, valign=center}
\newtcolorbox{tlumbox}{size=minimal, colback=white, colframe=red, arc=0mm, width=\linewidth, halign=left, valign=top}

\usepackage{Tabbing}
\usepackage{longtable}
\usepackage[savepos]{zref}

\let\grelocalleftbox\localleftbox
\let\grelocalrightbox\localrightbox

\usepackage[latin,polish]{babel}

\usepackage{fontspec}
\defaultfontfeatures{Ligatures=TeX}
\setmainfont{Libertinus Serif}

\usepackage{libertine} % ust. czcionki
\pagestyle{empty}	% usuwanie nagłówków/stopek z dokumentu
\usepackage[indent=0pt]{parskip} % ustawianie paragrafów
%\setstretch{1.5} 

\usepackage{graphicx} % umożliwia import obrazów
\graphicspath{{img/}} % katalog obrazów
\usepackage{pgfornament}% ornamenty graficzne

\usepackage{titlesec} % ustawianie nagłówków
\titlespacing*{\chapter}{0pt}{0pt}{0pt}	
\titleformat{\chapter}[hang]{\filcenter\LARGE\scshape\bfseries}{}{0pt}{} % format nagłówków
\titlespacing*{\section}{0pt}{0pt}{0mm}	
\titleformat{\section}[hang]{\filcenter\large\scshape\bfseries}{}{0pt}{} % format nagłówków

%\newfontface\vectis{VectisMonumentalW01-Regular}
%\def\greinitialformat#1{\fontsize{20}{20}\vectis #1}

% --- komendy własne ---
\newcommand{\red}[1]{\textcolor{red}{#1}}
\newcommand{\black}[1]{\textcolor{black}{#1}}

\newcommand{\ornament}[1]{
	\begin{center}
		\vspace{25mm}
		\pgfornament[width=0.6cm]{#1}
		\newpage
	\end{center}}

\begin{document}

\setlength{\parindent}{0pt}
\gresetbracerendering{font}

\pagenumbering{gobble} % nie jest wyprowadzany żaden numer strony, chociaż numer jest nadal resetowany. Odniesienia do tej strony również są puste.
\begin{spacing}{0.7}
\begin{center}
	\section*{Jesu Redemptor omnium}
	\small\textit{Liber Usualis, Solesmes, 1961, p.365}
	\vspace{5mm}
\end{center}	
\end{spacing}
\begin{paracol}{2}
	\ensurevspace{4cm}
	\begin{nutybox}\greannotation{Hymn.}\greannotation{1.} \gabcsnippet{(c4) JE(c)su(d) Red(eg)ém(gf)ptor(ede) ó(f')mni(e)um,(d.) (;) Quem(ce) lu(g')cis(h) an(h)t{<e>e</e>}(g) o(g)rí(h!iwj)gi(ivHG)nem,(h.) (:) Pa(hi)rem(g) pa(fg)tér(e)nae(dc) gló(d)ri(de)ae,(e.) (;) Pa(c)ter(d) su(eg)pré(gf)mus(ede) é(f')di(e)dit.(d.) (::) }
	\end{nutybox}
	\switchcolumn
	\begin{tlumbox}
	\vspace{3.5mm} 
	Jezu, Odkupicielu wszystkich, którego przed powstaniem światła (jako) równego chwale Ojcowskiej Ojciec zrodził najwyższy.
	\end{tlumbox}
  	\ensurevspace{4cm}
	\switchcolumn*
	\begin{nutybox}
	\gresetinitiallines{0}\gabcsnippet{(c4) <c>2.</c> Tu(c) lu(d)men(eg) et(gf) splen(ede)dor(f') Pa(e)tris,(d.) (;) Tu(ce) spes(g') per(h)én(h)nis(g) ó(h!iwj)mni(ivHG)um :(h.) (:) In(hi)tén(g)de(fg) quas(e) fun(dc)dunt(d) pre(de)ces(e.) (;) Tu(c)i(d) per(eg) or(gf)bem(ede) sér(f')vu(e)li.(d.) (::) }
	\end{nutybox}
	\switchcolumn
	\begin{tlumbox}
	\vspace{3.5mm} 
	Ty światłości i~odblasku Ojcowski, Ty nadziejo nieustanna wszystkich, wysłuchaj modłów, jakie słudzy twoi wylewają po (całym) świecie.
	\end{tlumbox}
	\ensurevspace{4cm}
	\switchcolumn*
	\begin{nutybox}
	\gresetinitiallines{0}\gabcsnippet{(c4) <c>3.</c> Me(c)mén(d)to,(eg) re(gf)rum(ede) Cón(f')di(e)tor,(d.) (;) No(ce)stri(g') quod(h) o(h)lim(g) cór(h!iwj)po(ivHG)ris,(h.) (:) Sac(hi)rá(g)t{<e>a</e>}(g) ab(fg) al(e)vo(dc) Vír(d)gi(de)nis(e.) (;) Na(c)scén(d)do,(eg) for(gf)mam(ede) súm(f')pse(e)ris.(d.) (::)}
	\end{nutybox}
	\switchcolumn
	\begin{tlumbox}
	\vspace{3.5mm} 
	Pomnij, Stwórco (wszech)rzeczy, żeś ongiś, rodząc się z~świętego Dziewicy żywota, przyjął postać ciała naszego.
	\end{tlumbox}
	\ensurevspace{4cm}
	\switchcolumn*
	\begin{nutybox}
	\gresetinitiallines{0}\gabcsnippet{(c4) <c>4.</c> Te(c)stá(d)tur(eg) hoc(gf) prae(ede)sens(f') di(e)es,(d.) (;) Cur(ce)rens(g') per(h) an(h)ni(g) cír(h!iwj)cu(ivHG)lum,(h.) (:) Quod(hi) so(g)lus(fg) e(e) si(dc)nu(d) Pa(de)tris(e.) (;) Mun(c)di(d) sa(eg)lus(gf) ad(ede)vé(f')ne(e)ris.(d.) (::)}
	\end{nutybox}
	\switchcolumn
	\begin{tlumbox}
	\vspace{3.5mm} 
	Świadczy o~tem dzień dzisiejszy, wracający w~biegu roku, żeś sam jeden z~łona Ojca przyszedł na zbawienie świata.
	\end{tlumbox}
	\ensurevspace{4cm}
	\switchcolumn*
	\begin{nutybox}
	\gresetinitiallines{0}\gabcsnippet{(c4) <c>5.</c> Hunc(c) a(d)stra,(eg) tel(gf)lus,(ede) aé(f')quo(e)ra,(d.) (;) Hunc(ce) o(g')mne(h) quod(h) cae(g)lo(h!iwj) sub(ivHG)est,(h.) (:) Sa(hi)lú(g)tis(fg) au(e)ctó(dc)rem(d) no(de)vae(e.) (;) No(c)vo(d) sa(eg)lú(gf)tat(ede) cán(f')ti(e)co.(d.) (::)}
	\end{nutybox}
	\switchcolumn
	\begin{tlumbox}
	\vspace{3.5mm} 
	Jego gwiazdy, ziemia, morza, Jego wszystko, co jest pod niebem, (jako) Sprawcę nowego zbawienia nową pozdrawia pieśnią.
	\end{tlumbox}
	\ensurevspace{4cm}
	\switchcolumn*
	\begin{nutybox}
	\gresetinitiallines{0}\gabcsnippet{(c4) <c>6.</c> Et(c) nos,(d) be(eg)á(gf)ta(ede) quos(f') sa(e)cri(d.) (;) Ri(ce)gá(g')vit(h) un(h)da(g) sán(h!iwj)gui(ivHG)nis,(h.) (:) Na(hi)tá(g)lis(fg) ob(e) di(dc)em(d) tu(de)i,(e.) (;) Hy(c)mni(d) tri(eg)bú(gf)tum(ede) sól(f')vi(e)mus.(d.) (::) }
	\end{nutybox}
	\switchcolumn
	\begin{tlumbox}
	\vspace{3.5mm} 
	I~my, których zrosił strumień krwi świętej, z~powodu dnia narodzenia twego składamy daninę pieśni.
	\end{tlumbox}
	\ensurevspace{4cm}
	\switchcolumn*
	\begin{nutybox}
	\gresetinitiallines{0}\gabcsnippet{(c4) <c>7.</c> Je(c)su,(d) ti(eg)bi(gf) sit(ede) gló(f')ri(e)a,(d.) (;) Qui(ce) na(g')tus(h) es(h) de(g) Vír(h!iwj)gi(ivHG)ne,(h.) (:) Cum(hi) Pa(g)tr{<e>e</e>}(g) et(fg) al(e)mo(dc) Spí(d)ri(de)tu,(e.) (;) In(c) sem(d)pi(eg)tér(gf)na(ede) saé(f')cu(e)la.(d.) (::)  A(ded)men.(cd..) (::)}
	\end{nutybox}
	\switchcolumn
	\begin{tlumbox}
	\vspace{3.5mm} 
	Jezu, Tobie niechaj będzie chwała, który narodziłeś się z~Dziewicy, oraz Ojcu i~dobroczynnemu Duchowi, po nieskończone wieki. Amen.\footnotemark
	\end{tlumbox}
\end{paracol}	
\footnotetext{tłum. Ks. Bronisław Gładysz, "Hymny brewiarza rzymskiego oraz patronału polskiego. Przekład i objaśnienia", Poznań 1933}
\end{document}